% =====================================================================
% Chapter 8 — Route B Results
% =====================================================================
\chapter{Route~B Results}
\label{ch:resultsB}

This chapter reports the results of \routeB{}. It presents the training
progress of the three main runs (\secref{sec:resultsB_training}), the
qualitative epoch-by-epoch evidence for the dummy and drum
(\secref{sec:resultsB_epochs}), the inference behaviour on a real \suop{}
chair segment (\secref{sec:resultsB_chair}), and an interpretation of what the
loss curves do and do not prove (\secref{sec:resultsB_interpretation}).

\section{Training progress}
\label{sec:resultsB_training}

\tabref{tab:routeB_training} reports the training and validation Chamfer loss
of the three principal runs. All runs used the RTX~4060 Laptop GPU, batch size
32--48, Adam at learning rate $10^{-3}$ and 6000 training pairs with a 10\%
hold-out.

\begin{table}[H]
  \centering
  \small
  \begin{adjustbox}{max width=\textwidth}
  \begin{tabular}{@{}lccccccc@{}}
    \toprule
    Run & Data & Epochs & CD@ep1 & CD@final & Val CD@ep1 & Val CD@final & Total time \\
    \midrule
    Dummy (mesh-cond.) & \texttt{dataset\_dummy\_suop} & 10 & 0.03993 & 0.00112 & 0.00459 & 0.00073 & 1705~s \\
    Drum & \texttt{dataset\_drum\_suop} & 10 & 0.01435 & 0.00178 & 0.00362 & 0.00051 & 1642~s \\
    Chair & procedural & 150 & 0.01657 & 0.00228 & 0.00159 & 0.00048 & 14{,}827~s \\
    \bottomrule
  \end{tabular}
  \end{adjustbox}
  \caption{Route~B training summary. Symmetric Chamfer loss $L =
  \mathrm{CD}(C,G)+\mathrm{CD}(F,G)$ on unit-normalised clouds. All three runs
  converge: the final training CD is 1--2 orders of magnitude below its
  epoch-1 value, and validation CD stays below training CD, indicating no
  memorisation at the sampled checkpoints.}
  \label{tab:routeB_training}
\end{table}

The dummy run (trained on \routeA{}-mesh-conditioned pairs) shows the fastest
relative improvement, dropping from 0.03993 to 0.00112 within ten epochs; the
drum converges from a lower starting point (0.01435) to 0.00178. The chair
run, trained for 150 epochs, reaches a training CD of $2.28\times10^{-3}$
while the validation CD plateaus around $5\times10^{-4}$ after epoch 45 ---
the typical signature of a model that has learned the coarse structure and is
only polishing fine detail.

\section{Epoch-by-epoch evidence}
\label{sec:resultsB_epochs}

The epoch contact sheets (\figref{fig:dummy_epochs}, \figref{fig:drum_epochs})
render, for every epoch, the same three panels: the sonar-like partial input,
the complete reference (ground truth) and the network's fine completion.
Reading a sheet left-to-right shows the input, the truth, and what the network
produces from the input alone --- a direct visual check that the completion is
actually derived from the partial rather than from a memorised shape.

\begin{figure}[H]
  \centering
  \includegraphics[width=0.88\textwidth]{fig_dummy_epochs.png}
  \caption{Dummy \routeB{} training progression, epochs 1--10 (mesh-conditioned
  run). Each row shows the same semantic triple: partial input, complete
  reference, network completion. The completion evolves from a rough blob
  (epoch 1) to a recognisable mannequin silhouette with limbs (epoch 10).}
  \label{fig:dummy_epochs}
\end{figure}

\begin{figure}[H]
  \centering
  \includegraphics[width=0.88\textwidth]{fig_drum_epochs.png}
  \caption{Drum \routeB{} training progression, epochs 1--10. The cylindrical
  barrel and end caps emerge early; later epochs refine the end-cap rims and
  reduce blur at the edges of the completion.}
  \label{fig:drum_epochs}
\end{figure}

\figref{fig:dummy_ep05_vs_10} compares epoch 5 and epoch 10 for the dummy
directly: the mid-training completion has the correct gross proportions but
softer limbs, while the epoch-10 completion is crisper and closer to the
reference silhouette.

\begin{figure}[H]
  \centering
  \includegraphics[width=0.98\textwidth]{fig_dummy_ep05_vs_10.png}
  \caption{Dummy completion at epoch 5 versus epoch 10 (top to bottom: input,
  reference, completion; left column epoch 5, right column epoch 10). The
  later checkpoint sharpens the shoulder/arm junctions and legs.}
  \label{fig:dummy_ep05_vs_10}
\end{figure}

\section{Inference on a real SUOP partial}
\label{sec:resultsB_chair}

The trained chair checkpoint (\texttt{pcn\_chair.pth}) was applied to a real
\suop{} chair segment --- a single unprocessed partial view, the deployment
scenario that \routeB{} was designed for. \figref{fig:chair_pcn_progression}
shows the input and the completions at epochs 1 and 15.

\begin{figure}[H]
  \centering
  \includegraphics[width=0.98\textwidth]{fig_chair_pcn_progression.png}
  \caption{Route~B chair inference on a real \suop{} partial: input segment
  (1{,}595 points, 0.729~m diagonal), PCN completion at epoch 1 (0.659~m
  diagonal) and at epoch 15 (0.713~m diagonal). The completion grows toward the
  input scale while adding inferred structure in the unobserved regions.}
  \label{fig:chair_pcn_progression}
\end{figure}

Two quantitative observations are worth recording:

\begin{enumerate}
  \item \textbf{Scale recovery.} The completed cloud's bounding-box diagonal
  moves from 0.659~m at epoch 1 to 0.713~m at epoch 15, converging toward the
  0.729~m input diagonal. Because training normalises to the unit sphere, the
  network must \emph{re-learn} the absolute scale of the object category; the
  convergence of the output diagonal toward the input diagonal is evidence
  that it does.
  \item \textbf{Completeness.} The completion emits a full 4096-point (fine)
  cloud from a 1{,}595-point partial, filling the far side, the underside and
  parts of the frame that the single view never measured. This is the defining
  behaviour of \routeB{} and the direct contrast with \routeA{}'s open,
  evidence-limited output.
\end{enumerate}

\section{Interpretation: what these numbers do and do not prove}
\label{sec:resultsB_interpretation}

The results establish the following:

\begin{enumerate}
  \item \textbf{The architecture trains.} All runs reduce Chamfer loss
  substantially and monotonically, and validation stays below training,
  showing the PCN with coarse-to-fine decoding and folding is trainable on
  this data.
  \item \textbf{The prior is category-shaped.} The epoch sheets show the
  completions take on the correct category anatomy (mannequin silhouette,
  barrel shape, chair frame), which only a learned prior can supply from a
  partial view.
  \item \textbf{Inference generalises to real sonar.} The real chair partial
  is completed at the correct scale, and the result is visually coherent.
\end{enumerate}

What the numbers do \emph{not} prove is physical accuracy of the inferred
geometry. The Chamfer loss is computed against \emph{synthetic} references, so
a low loss means ``consistent with the training distribution'', not
``correct against the real object''. A visually plausible completion can be a
hallucination --- the network may interpolate between training shapes rather
than recover the actual object. The dummy's mesh-conditioned training set
mitigates this by matching the sensor's geometry distribution, but the real
\suop{} objects have no external scan with which to verify the completions.
This threat to validity is addressed in \secref{ch:discussion} and the future
work of \secref{ch:conclusions}.
